% Body of \section{Experiments} (\label{app:experiments}).
%
% Paste between \label{app:experiments} and \section{Notation}, or keep this
% file next to the code and \input it -- every number below is produced by
% `python run.py {benchmark,ablation,sweep}` and every figure is a file written
% into results/, so editing a number here by hand makes it disagree with the
% plots it sits next to.
%
% Needs \usepackage{graphicx} and \usepackage{booktabs}, and, if the images
% are not at the project root, \graphicspath{{results/}} in the preamble.
% Two figures are used: benchmark.png and ablation.png.  The sweeps are tables.
%
% Uses one macro the draft does not define. Add to the preamble:
%   \newcommand{\Gb}{\mathbf{G}}

Every instance draws the columns of $\Xb\in\RR^{d\times K}$ and
$\Ub\in\RR^{d_u\times K}$ i.i.d.\ standard normal and normalises them to unit
$\ell_2$ norm, and draws $\thetab_\star$ i.i.d.\ standard normal, rescaled so
that $\max_{a\in[K]}|\zb_a^\top\thetab_\star|=1$.  Instances whose gap to the
second best arm falls below $\Delta_2=0.10$ are redrawn: a near-tie inflates
$T_\delta$ without changing what the comparison measures, and would let one
draw dominate the mean.  Rewards are Gaussian with $\sigma=0.1$, drawn from a
per-arm stream so that the $n$th pull of an arm is the same variate whatever
the other arms did in between; the two algorithms therefore run on the same
instances with the same noise, and a difference between them is never a
difference in the draw.  Unless stated otherwise $K=20$, $d=8$, $\delta=0.05$,
and each figure uses $200$ instances with seeds $20260823,\dots,20261022$.  A
run that has not stopped after $T_{\max}=2\times10^5$ pulls is censored and
reported separately rather than being counted as a success.

\paragraph{Scenarios.}
Three scenarios differ only in $d_u$ and in the conditioning of the draw.  In
\emph{fully observable} we set $d_u=0$, so nothing is hidden and both methods
are well specified; this is the control.  In \emph{partially observable} we set
$d_u=d$ and draw without further conditioning, so the latent shifts reverse the
ordering on some instances and not others.  In \emph{partially observable
(reversed)} we set $d_u=d$ and draw conditional on
$\argmax{a\in[K]}\;\xb_a^\top\thetab_\star^{(o)}\neq a_\star$, so that the
recorded covariates on their own favour a suboptimal arm.  The conditioning
does not make the instances easier: their median $\Delta_2$ is $0.204$ against
$0.225$ in the unconditioned scenario.

\paragraph{Baseline.}
We compare against the canonical observed-features-only construction: the
$G$-optimal design $\pib^G$ of \cref{eq:design} computed once over
$\{\xb_a\}_{a\in[K]}$ and rounded, ridge regression fitted in $\RR^d$, and an
arm eliminated once the leader beats it by more than twice a confidence width,
\[
  \beta_t := \sigma\sqrt{2\log(2K/\delta)}\,\sqrt{d/t},
  \qquad
  \text{eliminate } a \text{ once }
  \big(\xb_{\psi}-\xb_a\big)^\top\hat\thetab_t > 2\beta_t .
\]
Each piece matches the corresponding piece of \texttt{CoLF}, so that the two
algorithms differ in their feature set and in nothing else.  The $\sqrt{d/t}$
is Kiefer--Wolfowitz: under $\pib^G$ we have
$\max_{a\in[K]}\|\xb_a\|^2_{\Vb(\pib^G)^{-1}}=d$, so after $t$ pulls allocated
to that design every arm's prediction variance is at most $d/t$.  The constant
$\sigma\sqrt{2\log(2K/\delta)}$ is the sub-Gaussian tail with a union bound
over the $K$ arms, the same tail \texttt{CoLF} uses; a bounded-noise constant
would be twice as wide and cost four times the samples.  The factor two
compares two arms, exactly as \cref{eq:stop} spends $\beta$ on each of the two
arms it compares.

The width is calibrated rather than merely conservative.  On the fully
observable scenario this agent is well specified, so any error there is the
width's fault and not the model's; over $150$ instances it errs on $2.0\%$ at
$\Delta_2\ge0.30$ and on none at $\Delta_2\ge0.10$, and on $0.5\%$ of the $200$
fully observable benchmark instances --- all below $\delta$, and none pinned at
zero.  A width that were too narrow would show it here, and a
width that were far too wide would show it as a rate of exactly zero at every
gap.

\paragraph{What the implementation realises.}
The code follows \cref{alg:colf} as written, with three departures.  Each is a
correction, and each is pinned by a regression test.

\emph{(i) The pseudo-action is drawn uniformly, $p=1/K$.}  At that value
\cref{eq:pseudo_policy} is the uniform draw --- $\phi_a=p=1/K$ on the pulled arm
and $(1-p)/(K-1)=1/K$ on every other --- which is the propensity a round-robin
visit carries.  Then $1/\phi_{a,\tau}=K$ and \cref{eq:dr_reward} is correct
exactly as written.  Away from $1/K$ it is not: at the draft's $p=0.6$ the
weight is $1/p=1.667$, one sweep leaves $(K-1/p)\Gb=18.33\,\Gb$ of imputation
error behind instead of cancelling it, and a width that covers that residual
honestly overshoots the realised error by a factor of $40$ to $233$ over
$t\in[20,2000]$, against $1.00$ to $1.18$ at $p=1/K$.  The redraw of
\cref{eq:rho} is unaffected in kind: the cap grows to $\rho_t=350$ at $t=10^3$
and the expected number of draws to $K$, both free, since a pseudo-draw costs
no pull.

\emph{(ii) Selection and stopping divide by $\sqrt{\pi_a}$.}  Since
$\hat\xb_a^\top\mub_\star=\sqrt{\pi_a}\,\zb_a^\top\thetab_\star$ by
\cref{eq:augment}, comparing $\hat\xb_a^\top\hat\mub_t$ directly compares
design-weighted rewards rather than rewards.  Even with a perfect estimate this
names the wrong arm on $19.5\%$ of instances at $K=20$ and on $49.5\%$ at
$K=50$.  \cref{eq:stop} is applied to $\hat\xb_a^\top\hat\mub_t/\sqrt{\pi_a}$
for the same reason.

\emph{(iii) The width is formed per arm before the maximum.}  Because $\hat\Xb$
is square and invertible, $\hat\Xb^\top\Gb^{-1}\hat\Xb=\Ib_K$ and hence
$\mathrm{Var}(\hat Y_a)=\sigma^2/S$ for \emph{every} arm, with $S$ the number
of completed sweeps.  The reward-scale budget is therefore the same number for
all arms.  Taking a maximum in $\hat\xb$-coordinates and dividing by
$\sqrt{\pi_a}$ afterwards inflates each arm's budget by
$\sqrt{\pi_{\max}/\pi_a}$, which is $1.3\times$ at $K=10$ and $2.1\times$ at
$K=50$, and costs the square of that in samples.

One consequence is worth recording.  Because $\hat\Xb$ is square, at $w=K$ the
imputation residual is $\zero$ to machine precision --- measured
$\|K\Gb-K\sum_a\hat\xb_a\hat\xb_a^\top\|_2=0$ exactly --- and the predicted
mean of each arm then coincides with that arm's sample mean to
$3\times10^{-16}$.  Removing the augmentation breaks this:
$\hat\Xb^\top\Gb^{-1}\hat\Xb$ is the projection $\Pb_{\bar\Xb}$ rather than
$\Ib_K$, and the same comparison is off by $0.34$.  The augmentation is what
lets an arbitrary per-arm offset be represented at all, which is the quantity
\cref{thm:sample_complexity_lower_bound} says the other arms cannot supply.

At these settings the width of \cref{eq:stop} reduces to
$\beta_{N}(\delta')=\sigma\sqrt{2K\log(2K/\delta')/N}$ with $N=N(\Mcal_{\le t})$
counting matched rounds, not sweeps: every matched round contributes the whole
of $\Gb$ to the information matrix, so after $N$ of them each arm carries
variance $\sigma^2K/N$.  \texttt{CoLF} tests \cref{eq:stop} only at sweep
boundaries, so $T_\delta$ is always a multiple of $K$.

\paragraph{Benchmark.}
\cref{fig:benchmark} and \cref{tab:benchmark} report the three scenarios.
\texttt{CoLF} certifies the correct arm on all $600$ runs.  The baseline is
$\delta$-sound when nothing is hidden --- one error in $200$ --- and once part
of the reward is hidden it certifies a suboptimal arm on $30.0\%$ of instances,
rising to $35.5\%$ when the observed covariates are made to favour one.  These
are six and seven times $\delta$.

The baseline is also three to four times faster, and that gap is the price of
\cref{thm:sample_complexity_lower_bound}.  Its estimate is $d$-dimensional, so a sample from one arm
sharpens its prediction at every other arm and it need never visit them all;
that is what the linear structure buys, and where the structure holds --- the
fully observable scenario --- it is both faster and correct.  \texttt{CoLF}
cannot do this.  Each arm carries its own latent offset, which by
\cref{thm:sample_complexity_lower_bound} the samples of the other arms do not identify, so every arm must
be visited and the complexity carries $K$ where the baseline carries $d$.  The
factor of three to four we measure is the factor $K/d$ the setting demands.
When the structure does not hold, the same speed is exactly what makes the
baseline wrong: it stops on a width that measures its noise and not the part of
the reward it cannot express.

\begin{figure}[t]
  \centering
  \includegraphics[width=\textwidth]{benchmark.png}
  \caption{Sample complexity (top) and failure rates (bottom) over $200$
  instances per scenario at $\delta=0.05$.  Boxes give the interquartile range
  and the bar the median.  Neither outliers nor the mean are drawn: censored
  runs sit at $T_{\max}$, where they would stretch the axis over four decades
  and would pull a mean above the box while not being a mean at all, since a
  censored run's true stopping time is unknown.  They are counted in the bottom
  row instead.}
  \label{fig:benchmark}
\end{figure}

\begin{table}[t]
  \centering
  \caption{Benchmark over $200$ instances per scenario at $\delta=0.05$.  A run
  that has not stopped by $T_{\max}=2\times10^5$ is censored and counted
  separately; \cref{fig:benchmark} shows the full distribution of $T_\delta$.}
  \label{tab:benchmark}
  \small
  \begin{tabular}{llrrr}
    \toprule
    scenario & method & median $T_\delta$ & error rate & censored\\
    \midrule
    fully observable
      & \texttt{CoLF}          & $220$ & $0.000$ & $0.000$\\
      & \textsc{G-alloc.}      & $82$  & $0.005$ & $0.005$\\
    \addlinespace
    partially observable
      & \texttt{CoLF}          & $240$ & $0.000$ & $0.000$\\
      & \textsc{G-alloc.}      & $66$  & $0.300$ & $0.000$\\
    \addlinespace
    partially obs.\ (reversed)
      & \texttt{CoLF}          & $280$ & $0.000$ & $0.000$\\
      & \textsc{G-alloc.}      & $70$  & $0.355$ & $0.010$\\
    \bottomrule
  \end{tabular}
\end{table}

\cref{thm:sample_complexity_lower_bound} is visible in the same runs.  \textsc{G-allocation} stops with
$T_\delta<K$ on a small fraction of runs, and on those runs it is almost always
wrong: an arm it never sampled is an arm whose latent offset the samples of the
other arms do not identify, so its reward is undetermined however large its
gap.  \texttt{CoLF} never stops before a sweep is complete and so never
violates the bound.

\paragraph{Ablation.}
\cref{fig:ablation} removes one component at a time on a common round-robin
schedule, so that the curves compare estimators rather than schedules.
\emph{Without augmentation} keeps the mixed design, the round-robin sampling,
the doubly robust estimator and the stopping rule, and removes only the $K-d$
appended directions of \cref{eq:augmented_matrix}; $\hat\Xb$ is then $d\times K$
rather than square, $\hat\Xb^\top\Gb^{-1}\hat\Xb$ becomes the projection
$\Pb_{\bar\Xb}$ instead of $\Ib_K$, and the estimator can no longer reproduce a
reward lying outside $R(\bar\Xb)$.  \emph{With true feature} hands
\textsc{G-allocation} the full $\zb_a=[\xb_a;\ub_a]$, the ceiling on what
seeing the latent covariates would buy.  Run to their own stopping rules on
$150$ instances of the reversed scenario:

\begin{center}
\begin{tabular}{lrrr}
  \toprule
  & error rate & censored & median $T_\delta$\\
  \midrule
  \textsc{G-allocation}                      & $0.367$ & $0.007$ & $68$\\
  \texttt{CoLF} without augmentation         & $0.720$ & $0.047$ & $380$\\
  \texttt{CoLF}                              & $0.000$ & $0.000$ & $280$\\
  \textsc{G-allocation} with true feature    & $0.027$ & $0.007$ & $204$\\
  \bottomrule
\end{tabular}
\end{center}

Two things follow.  The doubly robust estimator does not carry the result on
its own: without the augmentation it is wrong on $72\%$ of instances, twice as
often as the baseline it shares a feature set with, because it commits to a
misspecified $d$-dimensional model under a width that presumes the model
correct and then spends five times as long confirming it.  And handing the
latent features to the baseline recovers $\delta$-soundness --- $2.7\%$, below
$\delta$ --- which locates the failure in the feature set rather than in the
allocation or the elimination rule.  Asking directly whether \texttt{CoLF}
recovers $\ub_a^\top\thetab_\star^{(u)}$ is not possible, since
\cref{eq:reparam_mu} absorbs the component of the latent contribution lying in
$R(\bar\Xb)$ into $\mub_\star^{(\parallel)}$ and $\thetab_\star$ is not
identifiable; the counterfactual is well posed, and it answers that seeing
$\ub$ would have fixed the baseline without making it faster than
\texttt{CoLF}.

\begin{figure}[t]
  \centering
  \includegraphics[width=\textwidth]{ablation.png}
  \caption{Estimation quality on a common round-robin schedule, $100$ instances
  of the partially observable scenario.  Colour is the method family and dash
  is the variant within it.  Left: the largest predicted-mean error across
  arms.  Right: that same error divided by the confidence width, so that on
  both panels the error is the numerator and lower is better; above one the
  agent is confident and wrong.  Bands are the 10th--90th percentile across
  instances.}
  \label{fig:ablation}
\end{figure}

\paragraph{$\delta$-soundness.}
A single confidence level cannot separate an algorithm that is $\delta$-sound
from one that merely happens to be accurate enough at that level.  We therefore
vary $d_u$ and $\delta$ one at a time, holding everything else fixed and
reusing the same seeds at every level, so that a difference between two rows is
a difference in that factor alone.

\cref{tab:du} varies how much of the reward is hidden.  At $d_u=0$ both methods
are exact, which is the control: nothing is hidden and the baseline is well
specified.  A \emph{single} hidden dimension is enough to break that, and the
rate then stays roughly constant --- with $\|\ub_a\|=1$ the per-arm shift is
$N(0,1)$ for every $d_u\ge1$, so $d_u$ changes only how correlated the shifts
are across arms, and what saturates is the chance that they reorder the arms
rather than their size.  \texttt{CoLF} is exact at every level, having formed
no estimate of $d_u$ and appended $K-d$ directions throughout, which is the
sense in which the augmentation introduces no tuning parameter.

\begin{table}[t]
  \centering
  \caption{Hidden dimensions, $150$ instances per level of the partially
  observable scenario at $\delta=0.05$.  Everything except $d_u$ is held fixed
  and the same seeds are used at every level.  The last column is the fraction
  of instances on which no observed-only linear rule can name the right arm
  however long it samples, and is reported to show that it, rather than the
  size of the shifts, is what saturates.}
  \label{tab:du}
  \small
  \begin{tabular}{rrrrrr}
    \toprule
    & \multicolumn{2}{c}{error rate} & \multicolumn{2}{c}{median $T_\delta$} & \\
    \cmidrule(lr){2-3}\cmidrule(lr){4-5}
    $d_u$ & \texttt{CoLF} & \textsc{G-alloc.} & \texttt{CoLF} & \textsc{G-alloc.}
      & unsolvable\\
    \midrule
    $0$  & $0.000$ & $0.000$ & $220$ & $79$ & $0.00$\\
    $1$  & $0.000$ & $0.187$ & $270$ & $80$ & $0.33$\\
    $2$  & $0.000$ & $0.220$ & $280$ & $68$ & $0.38$\\
    $4$  & $0.000$ & $0.247$ & $270$ & $60$ & $0.47$\\
    $8$  & $0.000$ & $0.327$ & $260$ & $66$ & $0.50$\\
    $16$ & $0.000$ & $0.213$ & $210$ & $47$ & $0.53$\\
    \bottomrule
  \end{tabular}
\end{table}

\cref{tab:delta} varies the confidence level.  The baseline's error rate does
not respond to it: the rate moves from $0.365$ to $0.280$ while $\delta$ falls
by a factor of $200$, so what was a $1.8$-fold violation becomes a $280$-fold
one.  Its sample complexity does grow, by a factor $2.6$ over the same range,
so tightening $\delta$ buys samples and nothing else --- exactly the failure
described in the introduction, where the learner's uncertainty shrinks with the
samples but the shift it omits does not.  \texttt{CoLF} certifies the correct
arm on all $1200$ runs; $0/200$ at each level places a one-sided $95\%$
Clopper--Pearson upper bound of $0.015$ on its rate, which establishes
$\delta$-soundness down to $\delta=0.05$ and is consistent with, but does not
by itself establish, soundness below it.

\begin{table}[t]
  \centering
  \caption{Confidence level, $200$ instances per level of the reversed
  scenario.  ``err.\ / $\delta$'' is the realised error rate divided by the
  level it was run at; a $\delta$-sound algorithm keeps it at or below one.
  $p$ is a one-sided binomial test of the realised error count against
  $\delta$.  \texttt{CoLF} makes no error at any level.}
  \label{tab:delta}
  \small
  \begin{tabular}{rrrrrrr}
    \toprule
    & \multicolumn{4}{c}{\textsc{G-allocation}} & \multicolumn{2}{c}{\texttt{CoLF}}\\
    \cmidrule(lr){2-5}\cmidrule(lr){6-7}
    $\delta$ & error & err.\ / $\delta$ & $p$ & median $T_\delta$
             & error & median $T_\delta$\\
    \midrule
    $0.2$   & $0.365$ & $1.8$   & $5\!\times\!10^{-8}$   & $60$  & $0.000$ & $240$\\
    $0.1$   & $0.355$ & $3.5$   & $3\!\times\!10^{-22}$  & $67$  & $0.000$ & $260$\\
    $0.05$  & $0.355$ & $7.1$   & $1\!\times\!10^{-40}$  & $70$  & $0.000$ & $280$\\
    $0.01$  & $0.340$ & $34.0$  & $8\!\times\!10^{-83}$  & $103$ & $0.000$ & $340$\\
    $0.005$ & $0.305$ & $61.0$  & $4\!\times\!10^{-89}$  & $118$ & $0.000$ & $380$\\
    $0.001$ & $0.280$ & $280.0$ & $2\!\times\!10^{-118}$ & $154$ & $0.000$ & $440$\\
    \bottomrule
  \end{tabular}
\end{table}

\paragraph{A wider width does not buy soundness back.}
The obvious repair is to be more conservative: keep the observed features and
eliminate at $c\,\beta_t$ for some $c>2$.  It does not work, and
\cref{tab:inflate} shows why.  Widening the rule does lower the error rate, but
it bottoms out at $2\%$ and never reaches $\delta$-soundness in any useful
regime, while the sample complexity grows as $c^2$: at $c=32$ the baseline
needs a median of $31\,477$ samples, a hundred times \texttt{CoLF}'s $280$, and
is still wrong twice as often.  The reason is that $\beta_t$ shrinks like
$1/\sqrt t$ while the part of the reward the observed features cannot express
does not shrink at all.  On these instances it is $0.44$ on the reward scale,
and $c\,\beta_t$ falls below it at $t\approx 10^4$ for $c=32$ just as it does at
$t\approx 10^2$ for $c=2$; the constant chooses when the rule starts firing on
the wrong arm, not whether.

\begin{table}[t]
  \centering
  \caption{Widening the elimination rule on the reversed scenario, $150$
  instances at $\delta=0.05$.  $c=2$ is the rule the analysis asks for, and the
  one used everywhere else in this section.}
  \label{tab:inflate}
  \small
  \begin{tabular}{rrrr}
    \toprule
    rule & error rate & censored & median $T_\delta$\\
    \midrule
    $2\beta_t$  & $0.367$ & $0.007$ & $68$\\
    $4\beta_t$  & $0.160$ & $0.013$ & $425$\\
    $8\beta_t$  & $0.047$ & $0.000$ & $2\,002$\\
    $16\beta_t$ & $0.020$ & $0.000$ & $7\,797$\\
    $32\beta_t$ & $0.020$ & $0.007$ & $31\,477$\\
    \midrule
    \texttt{CoLF} & $0.000$ & $0.000$ & $280$\\
    \bottomrule
  \end{tabular}
\end{table}
